% Source-derived chapter 08: Proof of the main theorem
% Original source: tates-thesis-de-rham-setting
\chapter{Proof of the main theorem}
\label{tates-thesis-de-rham-setting-chapter-08}
\source{tates-thesis-de-rham-setting}

\begin{remark}[Source section]
\label{tates-thesis-de-rham-setting-chapter-08-source}
\source{tates-thesis-de-rham-setting}
\source{tates-thesis-de-rham-setting}
This chapter reproduces the corresponding section of the retrieved
LaTeX source, including its definitions, statements, examples, and
proofs. It has not yet been translated into Lean.
\notready
\end{remark}

% The source section title was: Proof of the main theorem
\label{s:main}
\source{tates-thesis-de-rham-setting}

\subsection{Overview}

In this section we prove Theorem \ref{t:main}. The idea is simple given the
work we have done so far: bootstrap from the functors $\Delta_n$ constructed
earlier.


\subsection{Arcs}\label{ss:delta-arcs}
\source{tates-thesis-de-rham-setting}

We begin by constructing a strongly $\fL^+\bG_m$-equivariant functor:
%
\[
\Delta_{\infty}:D^!(\fL^+\bA^1) \to \IndCoh^*(\sY).
\]

\subsubsection{Notation}

In what follows, we let:
%
\[
p_n:\fL^+ \bA^1 \to \fL_n^+\bA^1
\]

\noindent denote the structural map.
For $m \geq n$, we similarly let: 
%
\[
p_{n,m}:\fL_m^+\bA^1 \to \fL_n^+\bA^1 
\]

\noindent denote the structural map.

\subsubsection{Overview}

First, let us unwind what it means to construct such a functor
$\Delta_{\infty}$.

Recall from \cite{dmod} that by definition we have: 
%
\[
D^!(\fL^+\bA^1) \coloneqq \colim_{n \geq 0} \, D(\fL_n^+\bA^1)  
\]

\noindent where the structural functors are the standard 
$!$-pullback functors of $D$-module theory. We let: 
%
\[
p_n^!: D(\fL_n^+\bA^1) \to 
D^!(\fL^+\bA^1)
\]

\noindent denote the structural functor.

The above colimit is a colimit of $D^*(\fL^+\bG_m)$-module categories.
Moreover, as it is indexed by $\bZ^{\geq 0}$, it suffices to 
construct functors:
%
\begin{equation}\label{eq:term-wise}
\source{tates-thesis-de-rham-setting}
D(\fL_n^+\bA^1) \to \IndCoh^*(\sY) \in 
D^*(\fL^+\bG_m)\mod 
\end{equation}

\noindent and commutative diagrams:
%
\begin{equation}\label{eq:delta-infty-comm}
\source{tates-thesis-de-rham-setting}
\vcenter{\xymatrix{
D(\fL_n^+\bA^1) \ar[drr]\ar[d]_{p_{n,n+1}^!} && \\ 
D(\fL_{n+1}^+\bA^1) \ar[rr] && \IndCoh^*(\sY) 
}}
\end{equation}

\noindent of $D^*(\fL^+\bG_m)$-module categories for each $n$.
(Here $D^*(\fL^+\bG_m)\mod $ acts on $\IndCoh^*(\sY)$ via
geometric class field theory, as in \S \ref{s:cft}.)

\subsubsection{}

For us, by definition, the functor \eqref{eq:term-wise}
is $\Delta_n$. It remains to construct the commutative 
diagrams \eqref{eq:delta-infty-comm} as above. 

In \S \ref{sss:diff-func}-\ref{sss:diff-compat}, we will formulate 
a compatibility for these commutative diagrams that 
characterizes them uniquely. We then turn to proving their existence.

\subsubsection{Preliminary constructions with differential operators}\label{sss:diff-func}
\source{tates-thesis-de-rham-setting}

We digress to some general constructions with $D$-modules, continuing
the discussion from \S \ref{ss:diff-1}. 

Suppose $\pi:X \to Y$ is a morphism between smooth varieties. In this
case, one has a canonical natural isomorphism:
%
\[
\Oblv^{\ell} \pi^! \simeq \pi^*\Oblv^{\ell}
\]

\noindent of functors:
%
\[
D(Y) \to \QCoh(X).
\]

\noindent Here $\pi^*$ is the quasi-coherent pullback functor. We refer
to \cite{crystals} for the construction.

The canonical map:
%
\[
\sO_Y \to \Oblv^{\ell}\ind^{\ell}(\sO_Y) = \Oblv^{\ell}(D_Y)
\]

\noindent then gives rise to a map:
%
\[
\sO_X \to \Oblv^{\ell}(D_Y) \simeq 
\Oblv^{\ell}\pi^!(D_Y)
\]

\noindent so by adjunction a map:
%
\begin{equation}\label{eq:can-x-y}
\source{tates-thesis-de-rham-setting}
\on{can}:D_X \to \pi^!(D_Y).
\end{equation}

In the special case where $\pi$ is smooth, the map $\on{can}$
is between objects concentrated in cohomological degree $-\dim X$, and
is an epimorphism in that abelian category.

Finally, we note that if an algebraic group $G$ acts on $X$ and $Y$, and the map
$\pi$ is $G$-equivariant, the map $\on{can}$ upgrades canonically to a map 
of $G$-weakly equivariant $D$-modules (by functoriality of the above constructions).  

\subsubsection{A compatibility}\label{sss:diff-compat}
\source{tates-thesis-de-rham-setting}

Let us now return to our setting. For each $n$, the above
provides canonical morphisms:
%
\[
\on{can}_n:
D_{\fL_{n+1}^+ \bA^1} \to p_{n,n+1}^! D_{\fL_n^+ \bA^1} \in 
D(\fL_{n+1}^+ \bA^1)^{\fL_{n+1}^+ \bG_m,w}.
\]

Suppose we are given a commutative diagram \eqref{eq:delta-infty-comm}.
We then obtain a morphism:
%
\begin{equation}\label{eq:delta-can-n}
\source{tates-thesis-de-rham-setting}
\begin{gathered}
i_{n+1,*}^{\IndCoh}(\sO_{\sZ^{\leq n+1}}) \eqqcolon 
\Delta_{n+1}(D_{\fL_{n+1}^+ \bA^1}) \xar{\on{\Delta_{n+1}(\on{can}_n)}}
\Delta_{n+1}(p_{n,n+1}^! D_{\fL_n^+ \bA^1}) 
\overset{\eqref{eq:delta-infty-comm}}{\simeq} \\
\Delta_n(D_{\fL_n^+ \bA^1}) 
\coloneqq
i_{n,*}^{\IndCoh}(\sO_{\sZ^{\leq n}}) \in \IndCoh^*(\sY)^{\fL^+\bG_m,w} \simeq 
\IndCoh^*(\sY_{\log}).
\end{gathered}
\end{equation}

\noindent (We have omitted terms $\lambda_{-,*}^{\IndCoh}$ to simplify 
the notation. We are also abusing notation in letting
e.g. $\Delta_n$ denote the induced functor on 
weakly equivariant categories.)

On the other hand, there is another such map coming from the embedding
$\delta_{n+1}:\sZ^{\leq n} \into \sZ^{\leq n+1}$ and the adjunction map:
%
\begin{equation}\label{eq:delta-delta}
\source{tates-thesis-de-rham-setting}
\sO_{\sZ^{\leq n+1}} \to \delta_{n+1,*}(\sO_{\sZ^{\leq n}}).
\end{equation}

We will show:

\begin{prop}
\source{tates-thesis-de-rham-setting}
\notready\label{p:delta-infty-comm}
\source{tates-thesis-de-rham-setting}

There exist unique commutative diagrams \eqref{eq:delta-infty-comm}
so that the resulting map \eqref{eq:delta-can-n} is \eqref{eq:delta-delta}.

\end{prop}

In this remainder of this subsection, we prove Proposition \ref{p:delta-infty-comm}.

\subsubsection{}\label{sss:hc-comm-diags}
\source{tates-thesis-de-rham-setting}

Let us make a preliminary observation in a general setting. Fix $G$
an affine algebraic group.

Suppose that $A$ is an associative algebra with a 
$G$-action and a Harish-Chandra datum as in \S \ref{ss:hc}, so 
$G$ acts strongly on $A\mod$.

Suppose we are given $\sC,\sD \in G\mod$ and functors:
%
\[
\begin{gathered}
F_1:A\mod \to \sC \\
F_2:\sC \to \sD \\
F_3:A \mod \to \sD.
\end{gathered}
\]

We let:
%
\[
\sF_1 \coloneqq F_1(A) \in \sC^{G,w}, \sF_3 \coloneqq F_3(A) \in \sD^{G,w}.
\]

Moreover, to remove discussion of higher homotopical considerations,
we assume:

\begin{itemize}

\item $A$ is classical.

\item $\ul{\End}_{\sC}(\sF_1)$ and $\ul{\End}_{\sD}(\sF_3)$ are classical.

\end{itemize}

Then, as in \S \ref{ss:hc}, the 
functor $F_1$ is completely encoded in the datum of compact
object $\sF_1 \in \sC^{G,w}$ with its action of $A^{op} \in \Alg(\Rep(G))$ by 
endomorphisms (the data should satisfy the property that the action of
$\fg$ through $i:\fg \to A$ coincides with the obstruction to strong equivariance).
The same applies for $F_3$.

Therefore, we see that (under the homotopically simplifying assumptions
above), the data of a commutative diagram:
%
\[
\xymatrix{
A\mod \ar[d]_{F_1} \ar[drr]^{F_3} && \\ 
\sC \ar[rr]^{F_2} && \sD 
}
\]

\noindent is equivalent to giving an isomorphism:
%
\[
F_2(\sF_1) \simeq \sF_3 \in \sD^{G,w}
\]

\noindent compatible with the actions of $A^{op} \in \Rep(G)$ on both sides.
(In particular, under these assumptions, it is equivalent to construct
a commutative diagram of strongly or weakly $G$-equivariant categories.)

\subsubsection{}

We apply this in our setting as follows. 

First, the actions on the source categories factor through the localization 
$D^*(\fL^+\bG_m) \to D(\fL_{n+1}^+\bG_m)$, and the functors
factor through the subcategory $\IndCoh^*(\sY^{\leq n+1}) \subset \IndCoh^*(\sY)$,
through which $D(\fL_n^+\bG_m)$ acts. So the question only concerns the
action of an affine algebraic group, not an affine group scheme. 

We then see
it suffices to show that there is a unique isomorphism:
%
\begin{equation}\label{eq:comm-iso}
\source{tates-thesis-de-rham-setting}
\Delta_{n+1}(p_{n,n+1}^! D_{\fL_n^+ \bA^1}) 
\simeq 
\Delta_n(D_{\fL_n^+ \bA^1}) 
\coloneqq
i_{n,*}^{\IndCoh}(\sO_{\sZ^{\leq n}})
\end{equation}

\noindent such that: 

\begin{itemize}

\item The two resulting maps from $i_{n+1,*}^{\IndCoh}(\sO_{\sZ^{\leq n+1}})$
(one being \eqref{eq:delta-can-n}, the other being the canonical map)
coincide.

\item The two resulting actions of $W_n^{op}$ on 
$\Av_!^{\Gr_{\bG_m}^{\leq n}}i_{n,*}^{\IndCoh}(\sO_{\sZ^{\leq n}})$ (coming from 
the tautological action of 
$W_n^{op}$ on $D_{\fL_n^+\bA^1} \in D(\fL_n^+\bA^1)$, and functoriality of
$\Delta_{n+1}p_{n,n+1}^!$ and $\Delta_n$ respectively) coincide
under the isomorphism.

\end{itemize}

As the canonical map:
%
\[
i_{n+1,*}^{\IndCoh}(\sO_{\sZ^{\leq n+1}}) \to 
i_{n,*}^{\IndCoh}(\sO_{\sZ^{\leq n}}) \in \Coh(\sY_{\log})^{\heart}
\]

\noindent is an epimorphism, we see that the first point above 
implies that such an isomorphism is unique if it exists. 

To see existence, it is convenient to fix a coordinate $t$ on the
disc. We use the
notation $\alpha$ notation of \S \ref{sss:ff-induction-strat}. 

We have an evident short exact sequence:
%
\[
0 \to D_{\fL_{n+1}^+ \bA^1}(-1) \xar{-\cdot \partial\alpha_n}
D_{\fL_{n+1}^+ \bA^1} \xar{\on{can}_n} 
p_{n,n+1}^! D_{\fL_n^+ \bA^1} \to 0 
\]

\noindent in the abelian category 
$D(\fL_{n+1}^+\bA^1)^{\fL_{n+1}^+ \bG_m,w,\heart}[n+1]$.
Here in the first term, the twist $(-1)$ indicates that we modify the
weakly equivariant structure by tensoring with the representation 
$k(-1) \in \Rep(\fL^+\bG_m)$ obtained by restriction along $\ev$ from the
standard representation of $\bG_m$.

By construction of $\Delta_{n+1}$, on weakly equivariant
categories, we have:
%
\[
\begin{gathered}
\Delta_{n+1}(D_{\fL_{n+1}^+ \bA^1}) \simeq i_{n+1,*}^{\IndCoh}(\sO_{\sZ^{\leq n+1}}) \\
\Delta_{n+1}(D_{\fL_{n+1}^+ \bA^1}(-1)) 
\simeq i_{n+1,*}^{\IndCoh}\iota_*^{\IndCoh}(\sO_{\sZ^{\leq n+1}}).
\end{gathered}
\]

\noindent Moreover, it is easy to see that $\Delta_{n+1}(-\cdot \partial_{\alpha_n})$
goes to the map induced from:
%
\[
\iota_*^{\IndCoh}(\sO_{\sZ^{\leq n+1}}) \xar{b_{-n}\cdot-} 
\sO_{\sZ^{\leq n+1}},
\]

\noindent i.e., the left arrow in \eqref{eq:fund-ses-2}. From 
\eqref{eq:fund-ses-2}, it follows that there is a unique
isomorphism \eqref{eq:comm-iso} compatible with the projection from 
$i_{n+1,*}^{\IndCoh}(\sO_{\sZ^{\leq n+1}})$. It is straightforward
to check that this isomorphism is compatible with the action of
$W_n^{op}$, concluding the argument.

\subsection{Compatibility with translation}

We now establish an additional equivariance property 
for the functor $\Delta_{\infty}$.

Specifically, recall the positive loop space $\fL^{pos} \bG_m \subset \fL \bG_m$ 
for $\bG_m$, as defined in \S \ref{ss:pos-gr}. 
This space is an ind-closed submonoid of $\fL \bG_m$ containing
$\fL^+\bG_m$, so we have 
a fully-faithful monoidal functor:
%
\[
D^*(\fL^{pos} \bG_m) \subset D^*(\fL \bG_m).
\]

By construction, $\fL^{pos} \bG_m$ acts on $\fL^+\bA^1$, so 
$D^*(\fL^{pos} \bG_m)$ acts on $D^!(\fL^+\bA^1)$.

On the other hand, we have the local class field theory isomorphism:
%
\[
D^*(\fL \bG_m) \simeq \QCoh(\LocSys_{\bG_m})
\]

\noindent defining a (fully faithful, symmetric) monoidal functor:
%
\[
D^*(\fL^{pos} \bG_m) \to \QCoh(\LocSys_{\bG_m})
\]

\noindent and so an action of $D^*(\fL^{pos} \bG_m)$ on 
$\IndCoh^*(\sY)$.

Our goal is to canonically upgrade the strong 
$\fL^+\bG_m$-equivariance of $\Delta_{\infty}$ to make $\Delta_{\infty}$
into a morphism of $D^*(\fL^{pos} \bG_m)$-module categories.

\subsubsection{}

Before proceeding, we wish to pin down this equivariance property uniquely.

Suppose $\sC$ is a $D^*(\fL^{pos} \bG_m)$-module category. 
We obtain a canonical functor $T:\sC^{\fL^+\bG_m} \to \sC^{\fL^+\bG_m}$
given by the action of (the $\delta$ $D$-module at) 
some (equivalently any) uniformizer $t \in \fL^{pos} \bG_m$.

Our uniqueness property will involve the corresponding functors $T$ on both
sides. Below, we collect preliminary observations regarding how the
functor $T$ interacts with the source and target of our functor. 

\subsubsection{} 

Suppose $\sC = D^!(\fL^+\bA^1)$. A choice of uniformizer $t$ defines
a closed embedding $\mu_t:\fL^+\bA^1 \into \fL^+ \bA^1$ given by multiplication
by $t$. The corresponding functor $T$ above is given by the left adjoint
to the pullback functor:
%
\[
\mu_t^!:D^!(\fL^+\bA^1) \to D^!(\fL^+ \bA^1).
\]

\noindent I.e., in the notation of \cite{dmod}, we have $T = \mu_{t,*,!-dR}$.

By adjunction, we observe that there is a canonical map:
%
\begin{equation}\label{eq:can-unit-geom}
\source{tates-thesis-de-rham-setting}
\on{can}:T(\omega_{\fL^+\bA^1}) = \mu_{t,*,!-dR}(\omega_{\fL^+\bA^1}) = 
\mu_{t,*,!-dR}\mu_t^!(\omega_{\fL^+\bA^1}) \to 
\omega_{\fL^+\bA^1}.
\end{equation}

\subsubsection{}

Now suppose $\sC = \IndCoh^*(\sY)$. Then we have:
%
\[
\begin{gathered}
\IndCoh^*(\sY)^{\fL^+\bG_m} \simeq \IndCoh^*(\sY)_{\fL^+\bG_m} \simeq \\
\QCoh(\LocSys_{\bG_m}^{\leq 0}) \underset{\QCoh(\LocSys_{\bG_m})}{\otimes}
\IndCoh^*(\sY) \subset 
\IndCoh^*(\sY \underset{\LocSys_{\bG_m}}{\times} \LocSys_{\bG_m}^{\leq 0}) = 
\IndCoh^*(\sY_{\log}^{\leq 0}).
\end{gathered}
\]

Under the above
identifications, the functor $T$ is given by tensoring with 
the \emph{dual}\footnote{Because the Contou-Carr\`ere pairing
is skew symmetric, one has to fix a convention in setting up geometric
class field theory. The sign compatibility was implicitly fixed in 
\S \ref{s:aj}. It is a straightforward check to trace the constructions
to see that our sign conventions force the appearance of the dual to the
standard representation; essentially, we are on the ``dual" side of the
Contou-Carr\`ere duality.}
to the standard representation of $\bB \bG_m$, i.e., tensoring with the
bundle $\sO_{\sY^{\leq 0}}(1)$ considered earlier.

Now observe that there is a canonical map:
%
\begin{equation}\label{eq:can-unit-spec}
\source{tates-thesis-de-rham-setting}
\on{can}:T(\sO_{\sZ^{\leq 0}}) = 
\sO_{\sZ^{\leq 0}}(1) \to \sO_{\sZ^{\leq 0}}
\end{equation}

\noindent from \eqref{eq:o(1)}.

\subsubsection{}

We can now formulate our results precisely. 

Observe that by construction, we have
an isomorphism:
%
\[
\Delta_{\infty}(\omega_{\fL^+\bA^1}) = \Delta_0(\omega_{\fL_0^+\bA^1}) \simeq
i_{0,*}^{\IndCoh}(\sO_{\sZ^{\leq 0}}) \in \IndCoh^*(\sY_{\log}^{\leq 0}).
\]

\noindent (Here we consider $\omega_{\fL^+\bA^1}$ as a $\fL^+\bG_m$-equivariant
object.)

\begin{prop}
\source{tates-thesis-de-rham-setting}
\notready

There is a unique structure of 
$D^*(\fL^{pos}\bG_m)$-equivariance on the functor $\Delta_{\infty}$
such that:

\begin{itemize}

\item The underlying $D^*(\fL^+\bG_m)$-equivariance is the one
constructed in \S \ref{ss:delta-arcs}.

\item The map:
%
\[
T(\sO_{\sZ^{\leq 0}}) \to \sO_{\sZ^{\leq 0}} \in \IndCoh^*(\sY_{\log}^{\leq 0})
\]

\noindent obtained by functoriality from:
%
\[
\Delta_{\infty}(\on{can}:T(\omega_{\fL^+\bA^1}) \xar{\eqref{eq:can-unit-geom}}  
\omega_{\fL^+\bA^1})
\]

\noindent coincides with \eqref{eq:can-unit-spec}.

\end{itemize}

\end{prop}

\begin{proof}

We begin with the existence. For this, it is convenient to fix a 
uniformizer $t$ once and for all. Note that this choice
induces a morphism $\LocSys_{\bG_m} \to \bB \bG_m$ using 
Proposition \ref{p:locsys}. We let $\sO_{\LocSys_{\bG_m}}(-1)$ and
$\sY(-1)$ denote
the corresponding line bundles, and generally use notation as in
\S \ref{sss:z-pic}. 

We obtain a splitting:
%
\[
D^*(\fL^{pos} \bG_m) \simeq D^*(\fL^+ \bG_m) \otimes D(\bZ^{\geq 0})
\]

\noindent of (symmetric) monoidal DG categories. Therefore, we find
that $D^*(\fL^{pos} \bG_m)$-module categories are equivalent to
$D^*(\fL^+\bG_m)$-module categories with an endofunctor.
Therefore, the existence amounts to showing that $\Delta_{\infty}$
intertwines $\mu_{t,*,!-dR}$ with tensoring with 
$\sO_{\sY}(1)$.

Unwinding the constructions, this amounts to constructing commutative diagrams
of strong $\fL^+\bG_m$-module categories:
%
\[
\xymatrix{
D(\fL_n^+\bA^1) \ar[rr]^{\mu_{t,*,dR}}\ar[d]^{\Delta_n} && D(\fL_{n+1}^+\bA^1) \ar[d]^{\Delta_{n+1}} \\
\IndCoh^*(\sY) \ar[rr]^{\sO_{\sY}(1) \otimes -} &&  \IndCoh^*(\sY)
}
\] 

\noindent compatibly in $n$ (using the $p_{-}^!$ functors and 
the commutativity data of Proposition \ref{p:delta-infty-comm});
here we have also let $\mu_t$ denote the induced map
$\fL_n^+\bA^1 \into \fL_{n+1}^+\bA^1$ given by multiplication with $t$.

Using the logic of \eqref{sss:hc-comm-diags} (as in 
the proof of Proposition \ref{p:delta-infty-comm}),
we see that this amounts to constructing
isomorphisms:
%
\[
\sO_{\sZ^{\leq n}}(1) \simeq \Delta_{n+1}(\mu_{t,*,dR} D_{\fL_n^+\bA^1}) \in 
\IndCoh^*(\sY_{\log}^{\leq n+1})
\]

\noindent satisfying some compatibilities (that do not involve higher
homotopy theory).

For this, we observe that there is a standard short exact sequence:
%
\[
0 \to D_{\fL_{n+1}^+\bA^1} \xar{-\cdot \alpha_0} D_{\fL_{n+1}^+\bA^1}(1) \xar{}
\mu_{t,*,dR} D_{\fL_n^+\bA^1}[1] \to 0 
\]

\noindent in the abelian category $D(\fL_{n+1}^+\bA^1)^{\fL_{n+1}\bG_m,w,\heart}[n+1]$
(with notation as before). 
The map $-\cdot \alpha_0$ maps under $\Delta_{n+1}$ to the canonical
map:
%
\[
\sO_{\sZ^{\leq n+1}} \to \iota_*^{\IndCoh}(\sO_{\sZ^{\leq n+1}}).
\]

Therefore, we obtain an isomorphism:
%
\[
\Delta_{n+1}(\mu_{t,*,dR} D_{\fL_n^+\bA^1}) \simeq 
\Coker(\sO_{\sZ^{\leq n+1}}[-1] \to \iota_*^{\IndCoh}(\sO_{\sZ^{\leq n+1}}[-1])) 
\simeq 
\Ker(\sO_{\sZ^{\leq n+1}} \to \iota_*^{\IndCoh}(\sO_{\sZ^{\leq n+1}})).
\]

\noindent The latter is canonically isomorphism to 
$\delta_{n+1,*}^{\IndCoh}\sO_{\sZ^{\leq n}}(1)$ by \eqref{eq:fund-ses}, 
as desired.

This isomorphism is easily seen to be compatible with the actions of $W_n^{op}$ 
and with varying $n$, proving existence.

For uniqueness, we observe that $\Delta_{\infty}$ is fully faithful
(cf. Lemma \ref{l:colim-ff} below), so the observation follows from the 
observation that:
%
\[
k \isom \End_{\TwoEnd_{\fL^+\bG_m\mod}(D^!(\fL^+\bA^1))}(\mu_{t,!,dR}) \in \Gpd
\]

\noindent (which e.g. follows by replacing functors with kernels
on the square $\fL^+\bA^1 \times \fL^+\bA^1$), 
meaning that any isomorphism we construct is unique up to scalars;
clearly the second compatibility in the proposition pins down this
scalar multiple uniquely.

\end{proof}

\subsection{Conclusion of the argument}

We now have the following precise form of our main theorem.

\begin{thm}
\source{tates-thesis-de-rham-setting}
\notready\label{t:main-precise}
\source{tates-thesis-de-rham-setting}

There is a unique equivalence:
%
\[
\Delta:D^!(\fL\bA^1) \isom  \IndCoh^*(\sY) 
\]

\noindent of $D^*(\fL\bG_m)$-module categories such that the induced morphism:
%
\[
D^!(\fL^+\bA^1) \subset D^!(\fL\bA^1) \to \IndCoh^*(\sY) 
\]

\noindent of $D^*(\fL^{pos}\bG_m)$-module categories coincides with the functor
$\Delta_{\infty}$ from above.

\end{thm}

\begin{proof}

There is a natural functor:
%
\[
D^*(\fL\bG_m) \underset{D^*(\fL^{pos}\bG_m)}{\otimes} D^!(\fL^+\bA^1) \to 
D^!(\fL\bA^1)
\]

\noindent that we observe is an equivalence. Indeed, this follows
by choosing a coordinate $t$ on the disc as before, splitting
$\fL\bG_m$ as $\fL^+\bG_m \times \bZ$, which then implies: 
%
\[
D^*(\fL\bG_m) \underset{D^*(\fL^{pos}\bG_m)}{\otimes} D^!(\fL^+\bA^1) \simeq 
\colim \, \big(D^!(\fL^+\bA^1) \xar{\mu_{t,*,!-dR}} D^!(\fL^+\bA^1) 
\xar{\mu_{t,*,!-dR}} \ldots \big).
\]

\noindent Using $t^{-n}$ to identify the $n$th term here with 
$t^{-n} \cdot \fL^+ \bA^1 \subset \fL \bA^1$ and applying the
definition from \cite{dmod}, we see that this 
colimit canonically identifies with $D^!(\fL\bA^1)$ compatibly with the
natural functor considered above.

Therefore, there is a unique functor:
%
\[
\Delta:D^!(\fL\bA^1) \isom  \IndCoh^*(\sY) 
\]

\noindent of $D^*(\fL\bG_m)$-module categories restricting to 
$\Delta_{\infty}$ (compatibly with the $D^*(\fL^{pos}\bG_m)$-equivariance).
It remains to show that $\Delta$ is an equivalence.

First, we note that $\Delta_{\infty}$ is fully faithful by 
Proposition \ref{p:ff} and Lemma \ref{l:colim-ff} below. Next,
we deduce that $\Delta$ itself is fully faithful by another
application of Lemma \ref{l:colim-ff}. 

Therefore, it remains to show that $\Delta$ is essentially surjective.
Clearly its essential image contains each object:
%
\[
\pi_{n,*}^{\IndCoh}(\sO_{\sZ^{\leq n}}(m))
\]

\noindent for all $n \geq 0$, $m \in \bZ$. Therefore, it suffices
to show that $\IndCoh^*(\sY)$ is generated under colimits by these objects.
But this follows immediately from the definitions and Corollary \ref{c:z-gen}.

\end{proof}

Above, we used the following simple categorical lemma.

\begin{lem}
\source{tates-thesis-de-rham-setting}
\notready\label{l:colim-ff}
\source{tates-thesis-de-rham-setting}

Suppose we are given a filtered diagram $i \mapsto \sC_i$ of compactly 
generated DG categories, and a functor:
%
\[
F:\sC \coloneqq \underset{i}{\colim} \, \sC_i \to \sD \in \DGCat_{cont}
\]

\noindent preserving compact objects. Suppose each structural functor:
%
\[
\sC_i \to \sC_j
\]

\noindent and each composition:
%
\[
\sC_i \to \sC \to \sD
\]

\noindent is fully faithful. Then $F$ is fully faithful.

\end{lem}



\bibliography{bibtex.bib}{}
\bibliographystyle{alphanum}
